% ---------------------------------------------------------------------------
% Replacement for section 1.1 of report.tex.
% Fixes: (i) GPD / GFF / convolution used interchangeably and never defined;
%        (ii) no formula stating what is actually parametrised;
%        (iii) the u-d phase present in the fit but absent from the equations;
%        (iv) the t-slope tying between flavours never stated.
% ---------------------------------------------------------------------------

\section{The model}

\subsection{What is parametrised}

It is worth being explicit about the level at which this model lives, because
three different objects are easily confused.

\begin{enumerate}
\item The \emph{generalised parton distributions} themselves,
      $X^q(x,\xi,t)$ for $X \in \{H_T, \bar E_T, \dots\}$ and flavour
      $q \in \{u,d\}$.  These depend on the loop momentum fraction $x$.
\item The \emph{convolution} of a GPD with the hard subprocess,
      obtained by integrating $x$ out.
\item The \emph{channel helicity amplitudes}, obtained from the convolutions
      by combining flavours with the photon couplings and the meson
      wavefunction.
\end{enumerate}

\textbf{This model never descends to level 1.}  No integral over $x$ is
performed anywhere, and no GPD is assumed.  What is parametrised is the
convolution, level 2.  That is the whole point of the approach: the data are
asked what they require of the convolutions, without committing to a GPD
parametrisation that produces them.

The convolution is
\begin{equation}
\langle X\rangle^q(\xi,t,Q^2) \;=\;
\int_{-1}^{1}\! dx \; X^q(x,\xi,t)\, K(x,\xi),
\qquad
K(x,\xi) \;\sim\; \frac{1}{x-\xi+i\epsilon},
\label{eq:conv}
\end{equation}
where $K$ is the hard-scattering kernel.  Its precise form is not needed
below, because \eqref{eq:conv} is never evaluated -- only parametrised.  What
matters is a structural consequence of the $i\epsilon$: the convolution is a
\emph{complex number},
\begin{equation}
\operatorname{Im}\langle X\rangle^q = \pi\, X^q(\xi,\xi,t),
\qquad
\operatorname{Re}\langle X\rangle^q = \mathcal{P}\!\!\int\! dx\, X^q K ,
\label{eq:reim}
\end{equation}
the imaginary part coming from the pole and the real part from the principal
value.  Since the $x$-shapes of $u$ and $d$ differ, $\langle X\rangle^u$ and
$\langle X\rangle^d$ cannot be assumed to share a phase.

\paragraph{The ansatz.}
Each convolution is parametrised by a modulus and a phase, separately for each
flavour:
\begin{align}
\bigl|\langle X\rangle^q\bigr| &= N_q \,
   \exp\!\big[(b_q + b'_q \ln x_B)\, t\big]\;
   \big(Q^2\big)^{n_Q^q/2},
\label{eq:mod}\\[2pt]
\arg\langle X\rangle^{d} - \arg\langle X\rangle^{u} &= \varphi_X .
\label{eq:phase}
\end{align}
The overall phase of a channel is unobservable, so $u$ is taken real and only
$d$ carries $\varphi_X$.  The parameters mean:

\begin{center}
\begin{tabular}{ll}
\toprule
$N$        & normalisation \\
$b$        & $t$-slope extrapolated to $x_B = 1$ \\
$b'$       & the slope at a given $x_B$ is $b + b'\ln x_B$, so $-b'$ plays the
             role of the Regge $\alpha'$ \\
$n_Q$      & power of $Q^2$ \\
$\varphi$  & phase of the $d$ convolution relative to $u$ \\
\bottomrule
\end{tabular}
\end{center}

\noindent
This is the same slope convention as the Goloskokov--Kroll fits.  An older
parameter convention used $L = \ln x_B - \ln 0.15$, so that $b$ was the slope
at $x_B = 0.15$; files in that convention must be converted
($b_{\rm new} = b + 1.897\,b'$), never copied raw.

Three blocks are parametrised this way: $\langle H_T\rangle$ and
$\langle \bar E_T\rangle$, which build the transverse amplitudes, and a
longitudinal block $\langle L\rangle$, which builds $T_{00}^{(-)}$.  The
labels $H_T$ and $\bar E_T$ name the GPD being convolved; the object that
carries the parameters is always the convolution, never the GPD itself.

\paragraph{Assumptions, stated plainly.}
Equations~\eqref{eq:mod} and~\eqref{eq:phase} embody two approximations that
should not be hidden.  A single exponential in $t$ for the modulus, and a
$t$-independent phase, are together equivalent to assuming that
$\operatorname{Re}\langle X\rangle$ and $\operatorname{Im}\langle X\rangle$
share a common $t$-slope: if they did not, the modulus
$\sqrt{\smash[b]{\operatorname{Re}^2 + \operatorname{Im}^2}}$ would not be a
single exponential and $\varphi = \arctan(\operatorname{Im}/\operatorname{Re})$
would depend on $t$.  Nothing forces that common slope; it is an assumption of
convenience, and relaxing it costs one parameter per block
($\varphi \to \varphi_0 + \varphi_1 t$, as already used for the phase of
$T_{00}^{(-)}$).

\subsection{From convolutions to channel amplitudes}

The observables are built from six channel helicity amplitudes.  In the
notation of the $\pi^0$ amplitude note,
\[
\begin{array}{llll}
T_{00}^{(+)} & \text{longitudinal, recoil non-flip} &\quad
T_{00}^{(-)} & \text{longitudinal, recoil flip}\\
T_{01}^{(+)} & \text{natural transverse, non-flip} &\quad
T_{01}^{(-)} & \text{natural transverse, flip}\\
U_{01}^{(+)} & \text{unnatural transverse, non-flip} &\quad
U_{01}^{(-)} & \text{unnatural transverse, flip}
\end{array}
\]
At the level used here $T_{01}^{(-)} = U_{01}^{(-)} = D/2$, real and equal,
which is the forward-kinematics relation; the remaining four carry the fitted
moduli and phases.

The photon couples to quark $q$ with charge $e_q$, and the meson wavefunction
selects which combination survives.  With
$\pi^0 = (u\bar u - d\bar d)/\sqrt2$ and
$\eta_8 = (u\bar u + d\bar d - 2s\bar s)/\sqrt6$, the coefficient of each
convolution is $\sum_q (\text{meson amplitude}) \times e_q$.  The strange term
drops out, for the reason given below.  Writing $X_q \equiv |\langle X\rangle^q|$, this gives
\begin{equation}
\langle X\rangle_{\pi^0 p} = \frac{2X_u + X_d\,e^{i\varphi_X}}{3\sqrt2},
\qquad
\langle X\rangle_{\eta p} = \frac{2X_u - X_d\,e^{i\varphi_X}}{3\sqrt6\,k},
\qquad
\langle X\rangle_{\pi^0 n} = \frac{X_u + 2X_d\,e^{i\varphi_X}}{3\sqrt2}.
\label{eq:flavour}
\end{equation}
A superscript on $\langle X\rangle$ is a flavour, a subscript is a channel:
$\langle X\rangle^u$ and $\langle X\rangle^d$ are the convolutions of
\eqref{eq:conv}, while $\langle X\rangle_{\rm ch}$ is the channel combination
built from them.  The relative sign differs between $\pi^0$ and $\eta$
because $e_d < 0$ meets
$-d\bar d$ in the $\pi^0$ and $+d\bar d$ in the $\eta$.  The neutron entry is
not new physics: it is the same meson factor with isospin applied to the
\emph{target}, $X_u^n = X_d^p$ and $X_d^n = X_u^p$.

\paragraph{Why there is no strange term, and no gluon term.}
Pseudoscalar meson production filters the $C$-odd, valence-like combination of
quark GPDs -- the hard kernel enters antisymmetrised in $x \to -x$, so quark
and antiquark contributions appear with opposite sign.  A symmetric sea,
$q_{\rm sea} = \bar q_{\rm sea}$, therefore cancels identically in that
combination.  Strangeness in the nucleon is pure sea, so
$\langle X\rangle^s = 0$ to the extent that $s = \bar s$; the approximation
fails only through the strangeness asymmetry $s - \bar s$, which is small.
This is why the $-2s\bar s$ piece of $\eta_8$ can be dropped in
\eqref{eq:flavour} without an additional parameter.

The same $C$-parity argument removes the gluon contribution altogether, which
is what makes $\pi^0$ and $\eta$ production a comparatively clean probe of the
valence sector: unlike vector mesons, there is no gluon GPD entering at this
order.  The singlet component of the $\eta$ survives only through the
$\eta$--$\eta'$ mixing of the decay constants, which is carried by $k_L$ in
\eqref{eq:kL}, not by an independent amplitude.

The combination in \eqref{eq:flavour} is a sum of \emph{amplitudes}; the
squaring happens only later, in \eqref{eq:gram}.  The phase $\varphi_X$ is
therefore physical and not removable: at $\varphi_X = 0$ the $\eta$ amplitude
can pass through an exact zero when $X_d = 2X_u$, whereas for
$\varphi_X \neq 0$
\begin{equation}
\bigl|2X_u - X_d e^{i\varphi_X}\bigr|^2
 = 4X_u^2 + X_d^2 - 4X_uX_d\cos\varphi_X
\label{eq:floor}
\end{equation}
has a floor of $4X_u^2\sin^2\varphi_X$.  Note also that this expression is
\emph{even} in $\varphi_X$: the unpolarised structure functions cannot
determine the sign of the phase, which is fixed only by the interference terms
$\sig{LT}$ and $\sig{LT'}$.

The factor $k$ is not flavour algebra but meson normalisation, and it enters
as a divisor, so $1/k$ is an enhancement of the $\eta$ channel.  In the
longitudinal twist-2 sector the meson enters through its decay constant; in
the two-mixing-angle scheme ($\theta_8 = -21.2^\circ$, $\theta_1 = -9.2^\circ$,
$f_8 = 1.26 f_\pi$, $f_1 = 1.17 f_\pi$),
\begin{equation}
k_L^{-1} = \Big(\cos\theta_8 - \sqrt2\,\tfrac{f_1}{f_8}\sin\theta_1\Big)
           \frac{f_8}{f_\pi} = 1.439
\qquad\Longrightarrow\qquad k_L = 0.695 .
\label{eq:kL}
\end{equation}
In the transverse twist-3 sector the meson enters through the chiral-odd
twist-3 distribution amplitude, normalised by
$\mu_M \sim m_M^2/(m_u+m_d)$, giving $k_T = 0.863$.
% TODO(VPK): k_T = 0.863 is currently a bare constant in the code with no
% derivation on record.  Give it the same explicit one-line formula as k_L,
% or quote the source.  It multiplies every eta transverse amplitude.

The near cancellation $2X_u - X_d$ in the $\eta$ channel is what makes that
channel sensitive to the flavour separation, and it is the reason $\eta$ data
carry weight far beyond their number.

\paragraph{The amplitudes themselves.}
Writing $\xi = \frac{x_B}{2-x_B}\left(1 + M_p^2/Q^2\right)$,
$t' = t - t_{\min}$ and $A = A(Q^2,x_B)$ for the phase-space factor,
\begin{equation}
D = \sqrt{2A\,(1-\xi^2)}\;\langle H_T\rangle_{\rm ch}, \qquad
U_{01}^{(+)} = \sqrt{A\,\kappa}\;\langle \bar E_T\rangle_{\rm ch}, \qquad
|T_{00}^{(-)}| = \sqrt{A\,\kappa}\;\langle L\rangle_{\rm ch},
\label{eq:amps}
\end{equation}
with $\kappa = -t'/(8M_p^2)$ and the subscript ``ch'' the channel combination
of \eqref{eq:flavour}.  Four further parameters unlock what would otherwise be
exact relations: $\rho_{CE}$ and $\phi_{CE}$ give $T_{01}^{(+)}$ a modulus and
phase relative to $U_{01}^{(+)}$ (breaking $C = E$), $\phi_w$ is the phase of
$U_{01}^{(+)}$, and $\rho_{nf}, \phi_{nf}$ give the longitudinal non-flip
amplitude $T_{00}^{(+)}$, which vanishes at leading order.  The phase of
$T_{00}^{(-)}$ is $\delta_0 + \delta_1 t$, common to both flavours -- note
that this one phase is allowed a $t$-dependence, unlike $\varphi_X$.

\subsection{What is free and what is tied}

The flavours are not treated symmetrically, and the pattern is a model
assumption worth stating.  For fit \texttt{\fittag}:

\begin{center}
\begin{tabular}{lcccc}
\toprule
block & slope $u$ & slope $d$ & $d/u$ ratio & phase $\varphi$ [rad] \\
      & \multicolumn{2}{c}{[GeV$^{-2}$, at $x_B = 0.25$]} & & \\
\midrule
$\langle H_T\rangle$       & 1.491 & 1.491 (tied to $u$) & $-0.296$ & $-1.261$ \\
$\langle \bar E_T\rangle$  & 1.110 & 5.720 (free)        & $0.935$  & $0.230$  \\
$\langle L\rangle$         & 1.780 & $= u$               & $-0.584$ & $0.512$  \\
\bottomrule
\end{tabular}
\end{center}

\noindent
Only $\langle \bar E_T\rangle$ is allowed an independent $t$-slope per
flavour, and it uses that freedom: the $d$ slope comes out five times the $u$
slope.  For $\langle H_T\rangle$ and for the longitudinal block the two
flavours share one slope and differ only in normalisation and phase.  Whether
$\langle H_T\rangle$ would also want a split slope is not tested here.
